% GENERATED FILE. Edit canonical lessons, then run npm run generate:books.
% canonical-source-hash: fnv1a64:16cca3a74981998c
% canonical-lessons: SA-C377-ativa, SA-C377-atyantam, SA-C377-tivrah, SA-C377-nitaram, SA-C377-atikramati

\chapter{Very and Extremely}
\label{ch:sa-very-and-extremely}
\hlchaptermodality{377}

\begin{chapteropening}
\textbf{By the end of this chapter:} \emph{I can say how strong something will be if it happens.}

Complete the last lesson of chapter 377: I can say how strong something will be if it happens.
\end{chapteropening}

\section[\texorpdfstring{\sk{अतीव} (atīva)}{अतीव (atīva)}]{\sk{अतीव} (atīva) — exceedingly; very}

\label{lesson:SA-C377-ativa}



\begin{quote}
Before the new one: say the Sanskrit for control, then the Sanskrit for takes refuge in.
\end{quote}

\subsection*{You'll want to know: \sk{अतीव}}
\textbf{\sk{अतीव}} — \emph{atīva} — \textquotedblleft{}exceedingly; very\textquotedblright{}.

\begin{quote}
\sk{यदि} \sk{भवान्} \sk{आगच्छति}, \sk{तर्हि} \sk{अहम्} \sk{अतीव} \sk{प्रसन्नः} \sk{भविष्यामि।}
\end{quote}

\begin{quote}
\emph{yadi bhavān āgacchati, tarhi aham atīva prasannaḥ bhaviṣyāmi.}
\end{quote}

\begin{quote}
If you come, I will be very pleased.
\end{quote}

\begin{cousinweb}
\textbf{\sk{अतीव}} (\emph{atīva}) MW gives \emph{atīva} \textquotedblleft{}exceedingly, very\textquotedblright{} and gives no derivation, so its origin is uncertain (MW).
\end{cousinweb}

\begin{grammarlens}[title={What you've built}]
One more word for saying what happens if.
\end{grammarlens}

\subsection*{Your turn}
\begin{itemize}
\raggedright
  \item \emph{Say it:} \emph{atīva}
  \item \emph{Say it:} \emph{atīva}, once more
  \item \emph{Say it:} \emph{yadi bhavān āgacchati, tarhi aham atīva prasannaḥ bhaviṣyāmi.} — If you come, I will be very pleased.
  \item \emph{Recall:} say the Sanskrit for control, then the Sanskrit for takes refuge in, then say \emph{atīva} again
\end{itemize}

\begin{tcolorbox}[breakable,colback=teal!4,colframe=teal!35!black,title={Before you move on}]
What does \sk{अतीव} mean? (\textquotedblleft{}Exceedingly; very\textquotedblright{}.) Say it once more.
\end{tcolorbox}
\section[\texorpdfstring{\sk{अत्यन्तम्} (atyantam)}{अत्यन्तम् (atyantam)}]{\sk{अत्यन्तम्} (atyantam) — extremely}

\label{lesson:SA-C377-atyantam}



\begin{quote}
Before the new one: say the Sanskrit for takes refuge in, then the Sanskrit for exceedingly.
\end{quote}

\subsection*{You'll want to know: \sk{अत्यन्तम्}}
\textbf{\sk{अत्यन्तम्}} — \emph{atyantam} — \textquotedblleft{}extremely\textquotedblright{}.

\begin{quote}
\sk{यदि} \sk{ग्रीष्मः} \sk{अत्यन्तम्} \sk{उष्णः}, \sk{तर्हि} \sk{जनाः} \sk{बहिः} \sk{न} \sk{गच्छन्ति।}
\end{quote}

\begin{quote}
\emph{yadi grīṣmaḥ atyantam uṣṇaḥ, tarhi janāḥ bahiḥ na gacchanti.}
\end{quote}

\begin{quote}
If the summer is extremely hot, people do not go out.
\end{quote}

\begin{cousinweb}
\textbf{\sk{अत्यन्तम्}} (\emph{atyantam}) is built from \emph{ati}, \textquotedblleft{}beyond\textquotedblright{}, and \emph{anta}, \textquotedblleft{}an end\textquotedblright{}; MW gives \emph{aty-anta} \textquotedblleft{}excessive, very great, very strong\textquotedblright{} and \textquotedblleft{}endless\textquotedblright{} (MW).
\end{cousinweb}

\begin{grammarlens}[title={What you've built}]
One more word for saying what happens if.
\end{grammarlens}

\subsection*{Your turn}
\begin{itemize}
\raggedright
  \item \emph{Say it:} \emph{atyantam}
  \item \emph{Say it:} \emph{atyantam}, once more
  \item \emph{Say it:} \emph{yadi grīṣmaḥ atyantam uṣṇaḥ, tarhi janāḥ bahiḥ na gacchanti.} — If the summer is extremely hot, people do not go out.
  \item \emph{Recall:} say the Sanskrit for takes refuge in, then the Sanskrit for exceedingly, then say \emph{atyantam} again
\end{itemize}

\begin{tcolorbox}[breakable,colback=teal!4,colframe=teal!35!black,title={Before you move on}]
What does \sk{अत्यन्तम्} mean? (\textquotedblleft{}Extremely\textquotedblright{}.) Say it once more.
\end{tcolorbox}
\section[\texorpdfstring{\sk{तीव्रः} (tīvraḥ)}{तीव्रः (tīvraḥ)}]{\sk{तीव्रः} (tīvraḥ) — intense; severe}

\label{lesson:SA-C377-tivrah}



\begin{quote}
Before the new one: say the Sanskrit for exceedingly, then the Sanskrit for extremely.
\end{quote}

\subsection*{You'll want to know: \sk{तीव्रः}}
\textbf{\sk{तीव्रः}} — \emph{tīvraḥ} — \textquotedblleft{}intense; severe\textquotedblright{}.

\begin{quote}
\sk{यदि} \sk{ज्वरः} \sk{तीव्रः}, \sk{तर्हि} \sk{वैद्यम्} \sk{आह्वयतु।}
\end{quote}

\begin{quote}
\emph{yadi jvaraḥ tīvraḥ, tarhi vaidyam āhvayatu.}
\end{quote}

\begin{quote}
If the fever is severe, call the doctor.
\end{quote}

\begin{cousinweb}
\textbf{\sk{तीव्रः}} (\emph{tīvraḥ}) MW gives \emph{tīvra} \textquotedblleft{}strong, severe, violent, intense\textquotedblright{} and takes it to come from the root √tu, \textquotedblleft{}to be strong\textquotedblright{} (MW).
\end{cousinweb}

\begin{grammarlens}[title={What you've built}]
One more word for saying what happens if.
\end{grammarlens}

\subsection*{Your turn}
\begin{itemize}
\raggedright
  \item \emph{Say it:} \emph{tīvraḥ}
  \item \emph{Say it:} \emph{tīvraḥ}, once more
  \item \emph{Say it:} \emph{yadi jvaraḥ tīvraḥ, tarhi vaidyam āhvayatu.} — If the fever is severe, call the doctor.
  \item \emph{Recall:} say the Sanskrit for exceedingly, then the Sanskrit for extremely, then say \emph{tīvraḥ} again
\end{itemize}

\begin{tcolorbox}[breakable,colback=teal!4,colframe=teal!35!black,title={Before you move on}]
What does \sk{तीव्रः} mean? (\textquotedblleft{}Intense; severe\textquotedblright{}.) Say it once more.
\end{tcolorbox}
\section[\texorpdfstring{\sk{नितराम्} (nitarām)}{नितराम् (nitarām)}]{\sk{नितराम्} (nitarām) — completely; wholly}

\label{lesson:SA-C377-nitaram}



\begin{quote}
Before the new one: say the Sanskrit for extremely, then the Sanskrit for intense.
\end{quote}

\subsection*{You'll want to know: \sk{नितराम्}}
\textbf{\sk{नितराम्}} — \emph{nitarām} — \textquotedblleft{}completely; wholly\textquotedblright{}.

\begin{quote}
\sk{यदि} \sk{भवान्} \sk{नितरां} \sk{श्रान्तः}, \sk{तर्हि} \sk{विश्रामं} \sk{करोतु।}
\end{quote}

\begin{quote}
\emph{yadi bhavān nitarāṃ śrāntaḥ, tarhi viśrāmaṃ karotu.}
\end{quote}

\begin{quote}
If you are completely worn out, take a rest.
\end{quote}

\begin{cousinweb}
\textbf{\sk{नितराम्}} (\emph{nitarām}) is built on \emph{ni}, \textquotedblleft{}down\textquotedblright{}, with the comparative ending \emph{-tarām}; MW gives \emph{ni-tarām} \textquotedblleft{}completely, wholly, entirely\textquotedblright{} (MW).
\end{cousinweb}

\begin{grammarlens}[title={What you've built}]
One more word for saying what happens if.
\end{grammarlens}

\subsection*{Your turn}
\begin{itemize}
\raggedright
  \item \emph{Say it:} \emph{nitarām}
  \item \emph{Say it:} \emph{nitarām}, once more
  \item \emph{Say it:} \emph{yadi bhavān nitarāṃ śrāntaḥ, tarhi viśrāmaṃ karotu.} — If you are completely worn out, take a rest.
  \item \emph{Recall:} say the Sanskrit for extremely, then the Sanskrit for intense, then say \emph{nitarām} again
\end{itemize}

\begin{tcolorbox}[breakable,colback=teal!4,colframe=teal!35!black,title={Before you move on}]
What does \sk{नितराम्} mean? (\textquotedblleft{}Completely; wholly\textquotedblright{}.) Say it once more.
\end{tcolorbox}
\section[\texorpdfstring{\sk{अतिक्रामति} (atikrāmati)}{अतिक्रामति (atikrāmati)}]{\sk{अतिक्रामति} (atikrāmati) — exceeds; goes beyond}

\label{lesson:SA-C377-atikramati}



\begin{quote}
Before the new one: say the Sanskrit for intense, then the Sanskrit for completely.
\end{quote}

\subsection*{You'll want to know: \sk{अतिक्रामति}}
\textbf{\sk{अतिक्रामति}} — \emph{atikrāmati} — \textquotedblleft{}exceeds; goes beyond\textquotedblright{}.

\begin{quote}
\sk{यदि} \sk{नदी} \sk{सीमाम्} \sk{अतिक्रामति}, \sk{तर्हि} \sk{ग्रामे} \sk{भयं} \sk{भविष्यति।}
\end{quote}

\begin{quote}
\emph{yadi nadī sīmām atikrāmati, tarhi grāme bhayaṃ bhaviṣyati.}
\end{quote}

\begin{quote}
If the river goes beyond its banks, there will be fear in the village.
\end{quote}

Past (it …): \sk{नदी} \sk{सीमाम्} \sk{अतिक्रान्तवती}, \emph{nadī sīmām atikrāntavatī}.

\begin{cousinweb}
\textbf{\sk{अतिक्रामति}} (\emph{atikrāmati}) is built on \emph{ati}, \textquotedblleft{}beyond\textquotedblright{}, and the root √kram, \textquotedblleft{}to step\textquotedblright{}; MW gives \emph{ati-√kram} \textquotedblleft{}to step or go beyond, to pass\textquotedblright{} and \emph{ati-krānta} \textquotedblleft{}having passed or transgressed\textquotedblright{} (MW).
\end{cousinweb}

\begin{grammarlens}[title={What you've built}]
One more word for saying what happens if.
\end{grammarlens}

\subsection*{Your turn}
\begin{itemize}
\raggedright
  \item \emph{Say it:} \emph{atikrāmati}
  \item \emph{Say it:} \emph{atikrāmati}, once more
  \item \emph{Say it:} \emph{yadi nadī sīmām atikrāmati, tarhi grāme bhayaṃ bhaviṣyati.} — If the river goes beyond its banks, there will be fear in the village.
  \item \emph{Recall:} say the Sanskrit for intense, then the Sanskrit for completely, then say \emph{atikrāmati} again
\end{itemize}

\begin{tcolorbox}[breakable,colback=teal!4,colframe=teal!35!black,title={Before you move on}]
What does \sk{अतिक्रामति} mean? (\textquotedblleft{}Exceeds; goes beyond\textquotedblright{}.) Say it once more.
\end{tcolorbox}
